\documentclass[11pt, a4paper]{article}

% bfw-style.tex — gemeinsame Style-Definitionen für alle Handouts/Aufgabenhefte
% Voraussetzung: Kompilation mit xelatex (wegen fontspec)

\usepackage[ngerman]{babel}
\usepackage{geometry}
\geometry{a4paper, top=2.2cm, bottom=2.2cm, left=2.3cm, right=2.3cm}

\usepackage{fontspec}
% Fallback-Kette: Source Sans Pro -> Calibri -> Segoe UI -> Latin Modern Sans
% (Latin Modern ist immer mit MiKTeX vorhanden)
\IfFontExistsTF{Source Sans Pro}{%
  \setmainfont{Source Sans Pro}[Ligatures=TeX]%
}{\IfFontExistsTF{Calibri}{%
  \setmainfont{Calibri}[Ligatures=TeX]%
}{\IfFontExistsTF{Segoe UI}{%
  \setmainfont{Segoe UI}[Ligatures=TeX]%
}{%
  \setmainfont{Latin Modern Sans}[Ligatures=TeX]%
}}}
\IfFontExistsTF{Source Code Pro}{%
  \setmonofont{Source Code Pro}[Scale=0.92]%
}{\IfFontExistsTF{Consolas}{%
  \setmonofont{Consolas}[Scale=0.92]%
}{%
  \setmonofont{Latin Modern Mono}[Scale=0.92]%
}}

\usepackage{xcolor}
% Farbpalette: hell, freundlich, modern
\definecolor{bfwPrimary}{HTML}{0EA5A4}    % Petrol/Türkis - Primär
\definecolor{bfwPrimaryDark}{HTML}{0F766E} % dunkler Petrol für Headings
\definecolor{bfwAccent}{HTML}{F59E0B}     % Amber - Akzent
\definecolor{bfwSuccess}{HTML}{10B981}    % Grün - Erfolg / Lösung
\definecolor{bfwWarn}{HTML}{F97316}       % Orange - Achtung
\definecolor{bfwInk}{HTML}{0F172A}        % fast Schwarz
\definecolor{bfwInkSoft}{HTML}{475569}    % gedämpftes Grau
\definecolor{bfwBgSoft}{HTML}{F8FAFC}     % off-white
\definecolor{bfwBgTeal}{HTML}{ECFEFF}     % helles Türkis
\definecolor{bfwBgAmber}{HTML}{FEF3C7}    % helles Gelb
\definecolor{bfwBgRose}{HTML}{FFE4E6}     % helles Rosé für Eselsbrücken

\usepackage[colorlinks=true, linkcolor=bfwPrimaryDark, urlcolor=bfwPrimaryDark]{hyperref}
\usepackage{enumitem}
\setlist[itemize]{leftmargin=*, itemsep=2pt, topsep=2pt}
\setlist[enumerate]{leftmargin=*, itemsep=2pt, topsep=2pt}

\usepackage[most]{tcolorbox}
\usepackage{booktabs}
\usepackage{tikz}
\usetikzlibrary{decorations.pathmorphing, positioning, shapes.geometric, arrows.meta}

% Merkkästchen — wichtige Definition / Kernpunkt
\newtcolorbox{merksatz}[1][]{%
  enhanced, breakable,
  colback=bfwBgTeal,
  colframe=bfwPrimary,
  boxrule=0.6pt,
  arc=4pt,
  left=10pt, right=10pt, top=8pt, bottom=8pt,
  fonttitle=\bfseries\color{bfwPrimaryDark},
  title={\faLightbulb\ Merksatz},
  #1
}

% Eselsbrücke — Mnemoniks
\newtcolorbox{eselsbruecke}[1][]{%
  enhanced, breakable,
  colback=bfwBgRose,
  colframe=bfwAccent,
  boxrule=0.6pt,
  arc=4pt,
  left=10pt, right=10pt, top=8pt, bottom=8pt,
  fonttitle=\bfseries\color{bfwWarn},
  title={Eselsbrücke},
  #1
}

% Beispiel-Box
\newtcolorbox{beispiel}[1][]{%
  enhanced, breakable,
  colback=bfwBgSoft,
  colframe=bfwInkSoft,
  boxrule=0.4pt,
  arc=3pt,
  left=10pt, right=10pt, top=8pt, bottom=8pt,
  fonttitle=\bfseries\color{bfwInk},
  title={Beispiel},
  #1
}

% Lösungs-Box (für Aufgabenhefte)
\newtcolorbox{loesung}[1][]{%
  enhanced, breakable,
  colback=bfwBgAmber!40,
  colframe=bfwSuccess,
  boxrule=0.6pt,
  arc=3pt,
  left=10pt, right=10pt, top=8pt, bottom=8pt,
  fonttitle=\bfseries\color{bfwSuccess!70!black},
  title={Lösung},
  #1
}

% Glossar-Eintrag
\newcommand{\glossareintrag}[2]{%
  \par\noindent\textbf{\color{bfwPrimaryDark}#1} \enspace #2\par\smallskip
}

% Section-Stil — etwas farbiger als Standard
\usepackage{titlesec}
\titleformat{\section}
  {\Large\bfseries\color{bfwPrimaryDark}}
  {\thesection}{0.6em}{}
\titleformat{\subsection}
  {\large\bfseries\color{bfwPrimary}}
  {\thesubsection}{0.5em}{}
\titleformat{\subsubsection}
  {\normalsize\bfseries\color{bfwInk}}
  {\thesubsubsection}{0.4em}{}

% TikZ-Stil "skizzenhaft" — etwas weicher als Rough.js, aber stabil
\tikzset{
  roughbox/.style = {
    draw=bfwPrimaryDark, thick, fill=bfwBgTeal,
    rounded corners=4pt, inner sep=6pt
  },
  roughaccent/.style = {
    draw=bfwAccent, thick, fill=bfwBgAmber,
    rounded corners=4pt, inner sep=6pt
  },
  roughline/.style = {
    draw=bfwInkSoft, thick, -{Stealth[length=2.5mm]}
  }
}

% Bullet-Symbole (bewusst kein FontAwesome — vermeidet Font-Installations-Probleme)
\newcommand{\faLightbulb}{$\bullet$}
\newcommand{\faBookmark}{$\bullet$}

% Header/Footer
\usepackage{fancyhdr}
\pagestyle{fancy}
\fancyhf{}
\renewcommand{\headrulewidth}{0pt}
\fancyfoot[L]{\small\color{bfwInkSoft}\thepage}
\fancyfoot[R]{\small\color{bfwInkSoft} BFW M-064 Beschaffung im Büromanagement}


\title{\vspace*{-1.5cm}\Huge\bfseries\color{bfwPrimaryDark}Aufgabenheft Session~3\\[0.4em]
  \large\color{bfwInkSoft}\textnormal{Das Unternehmen: Firma, Rechtsformen \& Organisation --- Übungen im IHK-Klausurformat}}
\author{\textsf{Lernfeld 1 --- Die eigene Rolle im Betrieb mitgestalten \textperiodcentered{} \small\copyright\ Can Siebert\,\textperiodcentered\,2026}}
\date{Selbstlernmaterial}

\begin{document}
\maketitle
\thispagestyle{empty}

\begin{merksatz}[title={\bfseries So nutzen Sie dieses Aufgabenheft}]
Bearbeiten Sie alle Aufgaben zunächst \emph{ohne} in die Lösungen zu schauen. Schreiben
Sie Ihre Antworten in vollständigen Sätzen --- die IHK-Prüfung verlangt formulierte
Begründungen, nicht bloße Stichworte. Beachten Sie die \emph{Operatoren} (nennen,
erläutern, beurteilen, begründen) --- sie geben den geforderten Antwort-Tiefegrad vor.
Erst danach öffnen Sie den Lösungsteil ab Seite~\pageref{loesungen} und vergleichen.
\end{merksatz}

\tableofcontents
\vspace{1em}
\hrule

\section{Aufgaben}

\subsection*{Aufgabe~1 --- Der Kaufmannsbegriff (10 Punkte)}

\noindent\textbf{\color{bfwAccent}YouTube-Vorbereitung:}\enspace \href{https://www.youtube.com/results?search_query=Kaufmann+HGB+Istkaufmann+Kannkaufmann+Formkaufmann}{Kaufmann nach HGB} (\url{https://www.youtube.com/results?search_query=Kaufmann+HGB+Istkaufmann+Kannkaufmann+Formkaufmann})

\medskip
Das HGB unterscheidet drei Grundtypen von Kaufleuten.

\begin{enumerate}[label=(\alph*)]
  \item \textbf{Erläutern} Sie den Unterschied zwischen Istkaufmann, Kannkaufmann und Formkaufmann. \hfill (3~P.)
  \item \textbf{Erklären} Sie den Unterschied zwischen deklaratorischer und konstitutiver Wirkung einer Eintragung. \hfill (3~P.)
  \item \textbf{Ordnen} Sie zu: e.K., GmbH, AG, OHG, eG --- welcher Kaufmannstyp liegt jeweils vor? \hfill (2~P.)
  \item \textbf{Beurteilen} Sie, warum ein Kleingewerbetreibender die freiwillige Eintragung gut überlegen sollte. \hfill (2~P.)
\end{enumerate}

\subsection*{Aufgabe~2 --- Firma und Firmengrundsätze (12 Punkte)}

\noindent\textbf{\color{bfwAccent}YouTube-Vorbereitung:}\enspace \href{https://www.youtube.com/results?search_query=Firma+HGB+Firmengrunds%C3%A4tze+Firmenwahrheit}{Firma und Firmengrundsätze} (\url{https://www.youtube.com/results?search_query=Firma+HGB+Firmengrunds\%C3\%A4tze+Firmenwahrheit})

\medskip
Die Firma ist der Name, unter dem ein Kaufmann seine Geschäfte betreibt (§ 17 HGB).

\begin{enumerate}[label=(\alph*)]
  \item \textbf{Nennen} Sie die vier Firmenarten mit je einem Beispiel. \hfill (4~P.)
  \item \textbf{Erläutern} Sie die Firmengrundsätze Firmenwahrheit, Firmenausschließlichkeit und Firmenbeständigkeit. \hfill (4~P.)
  \item \textbf{Beurteilen} Sie: Darf eine Möbelfabrik ihre Firma in „Rheinland Reisebüro GmbH" ändern? \hfill (2~P.)
  \item \textbf{Erklären} Sie den Unterschied zwischen Firmenkern und Firmenzusatz. \hfill (2~P.)
\end{enumerate}

\subsection*{Aufgabe~3 --- Handelsregister und Publizität (12 Punkte)}

\noindent\textbf{\color{bfwAccent}YouTube-Vorbereitung:}\enspace \href{https://www.youtube.com/results?search_query=Handelsregister+Abteilung+A+B+Publizit%C3%A4t}{Handelsregister und § 15 HGB} (\url{https://www.youtube.com/results?search_query=Handelsregister+Abteilung+A+B+Publizit\%C3\%A4t})

\medskip
Das Handelsregister genießt öffentlichen Glauben (§ 15 HGB).

\begin{enumerate}[label=(\alph*)]
  \item \textbf{Ordnen} Sie e.K., OHG, KG, GmbH und AG der Abteilung A oder B zu. \hfill (3~P.)
  \item \textbf{Erläutern} Sie die Publizität des Handelsregisters nach § 15 HGB. \hfill (3~P.)
  \item \textbf{Beurteilen} Sie folgenden Fall: Der Widerruf einer Prokura ist eingetragen und bekanntgemacht; die frühere Prokuristin bestellt dennoch mit alter Vollmacht Waren. Ist das Unternehmen gebunden? \hfill (4~P.)
  \item \textbf{Nennen} Sie, in welchem Register Genossenschaften und Vereine eingetragen werden. \hfill (2~P.)
\end{enumerate}

\subsection*{Aufgabe~4 --- Rechtsform begründet wählen (16 Punkte)}

\noindent\textbf{\color{bfwAccent}YouTube-Vorbereitung:}\enspace \href{https://www.youtube.com/results?search_query=Rechtsformen+Vergleich+GmbH+OHG+KG+Haftung}{Rechtsformen im Vergleich} (\url{https://www.youtube.com/results?search_query=Rechtsformen+Vergleich+GmbH+OHG+KG+Haftung})

\medskip
Drei Gründerinnen wollen ein Start-up für Bürosoftware aufbauen. Sie wollen ihr
Privatvermögen schützen; das Startkapital ist zunächst knapp.

\begin{enumerate}[label=(\alph*)]
  \item \textbf{Vergleichen} Sie OHG, GmbH und UG (haftungsbeschränkt) hinsichtlich Haftung und Mindestkapital. \hfill (6~P.)
  \item \textbf{Begründen} Sie, welche Rechtsform Sie den Gründerinnen empfehlen. \hfill (4~P.)
  \item \textbf{Erläutern} Sie die Organe einer GmbH und einer AG. \hfill (4~P.)
  \item \textbf{Erklären} Sie den Unterschied in der Haftung von Komplementär und Kommanditist in der KG. \hfill (2~P.)
\end{enumerate}

\subsection*{Aufgabe~5 --- Prokura und Handlungsvollmacht (12 Punkte)}

\noindent\textbf{\color{bfwAccent}YouTube-Vorbereitung:}\enspace \href{https://www.youtube.com/results?search_query=Prokura+Handlungsvollmacht+Unterschied+HGB}{Prokura vs.\ Handlungsvollmacht} (\url{https://www.youtube.com/results?search_query=Prokura+Handlungsvollmacht+Unterschied+HGB})

\medskip
Vollmachten regeln, wer ein Unternehmen nach außen vertreten darf.

\begin{enumerate}[label=(\alph*)]
  \item \textbf{Erläutern} Sie Umfang, Erteilung und Eintragung der Prokura. \hfill (3~P.)
  \item \textbf{Erläutern} Sie die drei Arten der Handlungsvollmacht. \hfill (3~P.)
  \item \textbf{Nennen} Sie die jeweiligen Unterschriftszusätze (Prokura, Generalvollmacht, Art-/Einzelvollmacht). \hfill (3~P.)
  \item \textbf{Beurteilen} Sie: Welche Vollmacht besitzt ein Auszubildender, der regelmäßig kassiert? \hfill (3~P.)
\end{enumerate}

\subsection*{Aufgabe~6 --- Aufbau- und Ablauforganisation (13 Punkte)}

\noindent\textbf{\color{bfwAccent}YouTube-Vorbereitung:}\enspace \href{https://www.youtube.com/results?search_query=Aufbauorganisation+Instanz+Stabstelle+Leitungssystem}{Aufbauorganisation und Leitungssysteme} (\url{https://www.youtube.com/results?search_query=Aufbauorganisation+Instanz+Stabstelle+Leitungssystem})

\medskip
\begin{enumerate}[label=(\alph*)]
  \item \textbf{Erklären} Sie die Begriffe Stelle, Instanz und Stabstelle. \hfill (3~P.)
  \item \textbf{Erläutern} Sie den Unterschied zwischen Einlinien-, Stablinien- und Matrixsystem mit je einem Vor- und Nachteil. \hfill (4~P.)
  \item \textbf{Grenzen} Sie Aufbau- und Ablauforganisation voneinander \textbf{ab}. \hfill (3~P.)
  \item \textbf{Ordnen} Sie zu: Produktion, Buchhaltung, Kontrolle --- Kern-, Unterstützungs- oder Managementprozess? \hfill (3~P.)
\end{enumerate}

\vspace{1em}
\noindent\textbf{Punkte gesamt: 75}

\newpage

\section{Lösungen}\label{loesungen}

\subsection*{Lösung Aufgabe~1}
\begin{loesung}
\textbf{(a)} \emph{Istkaufmann} (§ 1 HGB): betreibt ein Handelsgewerbe, das einen
kaufmännisch eingerichteten Geschäftsbetrieb erfordert --- Kaufmann kraft Gesetzes.
\emph{Kannkaufmann} (§§ 2, 3 HGB): Kleingewerbetreibende und land-/forstwirtschaftliche
Unternehmen, die sich freiwillig eintragen lassen können. \emph{Formkaufmann}:
Kapitalgesellschaften, die kraft Rechtsform Kaufmann sind.

\textbf{(b)} \emph{Deklaratorisch} = rechtsbekundend: die Eintragung bestätigt nur eine
schon bestehende Rechtslage (z.\,B.\ Istkaufmann, Prokura). \emph{Konstitutiv} =
rechtsbegründend: die Rechtslage entsteht erst mit der Eintragung (z.\,B.\ GmbH, AG).

\textbf{(c)} e.K.\ = Istkaufmann; OHG = Istkaufmann (Personengesellschaft); GmbH, AG, eG
= Formkaufmann.

\textbf{(d)} Mit der Eintragung treffen ihn die Buchführungspflichten nach HGB sowie
besondere Untersuchungs- und Rügepflichten beim Kaufvertrag --- also mehr Pflichten und
Aufwand.
\end{loesung}

\subsection*{Lösung Aufgabe~2}
\begin{loesung}
\textbf{(a)} Personenfirma (z.\,B.\ Gilles e.K.); Sachfirma (z.\,B.\ Balthasar Transporte
GmbH); Mischfirma (z.\,B.\ Duisdorfer BüroKonzept KG); Fantasiefirma (z.\,B.\ Kalyptus OHG).

\textbf{(b)} \emph{Firmenwahrheit} (§ 18 HGB): keine irreführenden Angaben, Unterscheidungs\-kraft.
\emph{Firmenausschließlichkeit} (§ 30 HGB): jede neue Firma muss sich von bestehenden am
selben Ort deutlich unterscheiden. \emph{Firmenbeständigkeit} (§§ 22, 23 HGB): die Firma
darf beim Erwerb des Geschäfts fortgeführt, aber nicht ohne das Handelsgeschäft veräußert
werden.

\textbf{(c)} Nein --- das verstößt gegen die Firmenwahrheit: „Reisebüro" entspricht nicht
dem tatsächlichen Gegenstand (Möbelherstellung) und führt in die Irre.

\textbf{(d)} Der \emph{Firmenkern} enthält Namen, Gegenstand oder Fantasiebezeichnung; der
\emph{Firmenzusatz} erläutert das Gesellschaftsverhältnis --- vor allem den Rechtsform\-zusatz
(e.K., OHG, GmbH usw.).
\end{loesung}

\subsection*{Lösung Aufgabe~3}
\begin{loesung}
\textbf{(a)} Abteilung A: e.K., OHG, KG (Einzelkaufleute und Personengesellschaften).
Abteilung B: GmbH, AG (Kapitalgesellschaften).

\textbf{(b)} § 15 HGB: Solange eine eintragungspflichtige Tatsache nicht eingetragen und
bekanntgemacht ist, kann sie einem gutgläubigen Dritten nicht entgegengehalten werden.
Ist sie eingetragen und bekanntgemacht, muss ein Dritter sie gegen sich gelten lassen.
Jeder darf auf den Registerinhalt vertrauen.

\textbf{(c)} Nein, das Unternehmen ist nicht gebunden. Da der Widerruf ordnungsgemäß
eingetragen und bekanntgemacht wurde, muss der Geschäftspartner ihn gegen sich gelten
lassen; er muss sich direkt an die frühere Prokuristin halten.

\textbf{(d)} Genossenschaften ins Genossenschaftsregister, Vereine ins Vereinsregister.
\end{loesung}

\subsection*{Lösung Aufgabe~4}
\begin{loesung}
\textbf{(a)} \emph{OHG}: mind.\ 2 Personen, keine Mindesteinlage vorgeschrieben, aber
unbeschränkte, unmittelbare und solidarische Haftung mit dem Privatvermögen.
\emph{GmbH}: Mindeststammkapital 25.000 €, Haftung auf das Gesellschaftsvermögen
beschränkt. \emph{UG (haftungsbeschränkt)}: gründbar ab 1 € Stammkapital, ebenfalls
beschränkte Haftung, muss Gewinne bis 25.000 € zurücklegen.

\textbf{(b)} Empfehlung: \textbf{UG (haftungsbeschränkt)} --- sie schützt das Privat\-vermögen
(beschränkte Haftung) und ist trotz knappem Startkapital gründbar; später Umwandlung in
eine GmbH möglich. Wer sofort mehr Kapital/Seriosität will, wählt die GmbH.

\textbf{(c)} \emph{GmbH}: Gesellschafterversammlung (bestellt die Geschäftsführer) und
Geschäftsführer. \emph{AG}: Hauptversammlung (Aktionäre), Vorstand (Geschäftsführung),
Aufsichtsrat (Kontrolle).

\textbf{(d)} Der Komplementär haftet unbeschränkt mit dem Privatvermögen; der Kommanditist
haftet nur bis zur Höhe seiner Kapitaleinlage.
\end{loesung}

\subsection*{Lösung Aufgabe~5}
\begin{loesung}
\textbf{(a)} Die \emph{Prokura} (§§ 48 ff.\ HGB) wird nur vom Inhaber ausdrücklich erteilt,
ins Handelsregister eingetragen und ermächtigt zu allen gerichtlichen und außer\-gerichtlichen
Geschäften des Handelsgewerbes (außer Grundstücksgeschäften ohne besondere Befugnis).

\textbf{(b)} \emph{Generalvollmacht}: alle gewöhnlichen Geschäfte des Betriebs.
\emph{Art-/Gattungsvollmacht}: eine bestimmte Art von Geschäften (z.\,B.\ Kassieren).
\emph{Einzel-/Spezialvollmacht}: ein einzelnes Geschäft (z.\,B.\ einmaliger Einkauf).

\textbf{(c)} Prokura: „ppa." (per procura); Generalvollmacht: „i.\,V." (in Vollmacht);
Art-/Einzelvollmacht: „i.\,A." (im Auftrag).

\textbf{(d)} Ein Auszubildender, der regelmäßig kassiert, besitzt eine \emph{Artvollmacht}
(Handlungsvollmacht für eine bestimmte Art von Geschäften). Er unterschreibt mit „i.\,A.".
\end{loesung}

\subsection*{Lösung Aufgabe~6}
\begin{loesung}
\textbf{(a)} \emph{Stelle}: kleinste organisatorische Einheit, Aufgabenbündel für eine
Person. \emph{Instanz}: Stelle mit Leitungs- und Weisungsbefugnis. \emph{Stabstelle}:
beratende Stelle ohne Weisungsbefugnis.

\textbf{(b)} \emph{Einliniensystem}: nur ein Vorgesetzter je Stelle --- klar, aber lange
Dienstwege. \emph{Stabliniensystem}: Einliniensystem plus Stabstellen --- entlastet die
Leitung, aber möglicher Stab-Linie-Konflikt. \emph{Matrixsystem}: doppelte Unterstellung
nach Funktion und Objekt/Projekt --- flexibel, aber Kompetenzkonflikte durch
Doppelunterstellung.

\textbf{(c)} Die \emph{Aufbauorganisation} legt das Gerüst fest (Stellen, Instanzen,
Abteilungen, Organigramm). Die \emph{Ablauforganisation} regelt die zeitlich-räumlichen
Arbeitsabläufe (wer erledigt was wann und wo).

\textbf{(d)} Produktion = Kernprozess; Buchhaltung = Unterstützungsprozess; Kontrolle =
Managementprozess.
\end{loesung}

\end{document}
